\documentclass[../DoAn.tex]{subfiles}
\begin{document}

\begin{center}\bfseries\large ABSTRACT\end{center}
\vspace{0.5cm}
RoomHub supports instructors in finding and booking classrooms or meeting rooms, while administrators review requests and manage schedules. The project uses a spreadsheet-based scheduling process as a reference scenario and implements room search, booking, approval, notification and access-permission workflows. Test data model an eight-storey building, with Building C7 at Hanoi University of Science and Technology serving as a contextual reference. The floor plan and room records are created by the team; they do not represent the building's actual layout or a campus-wide deployment.

This pre-final report is prepared for course ET4710, Mobile Application Programming. The system comprises an Android application built with React Native and TypeScript, a Java 17/Spring Boot service exposing REST APIs, and an H2 database accessed through JPA. The current configuration selects server mode, in which the application exchanges data through the API. On-device storage remains available as a separate mode; changes made while offline are not automatically reconciled with server data. The public tunnel is intended for testing and does not establish production availability.

For rooms equipped with a smart lock, an administrator or an authorised user may create a seven-digit temporary PIN. The application records the PIN and its validity interval before attempting delivery. Once a OneIoT session is connected, the application publishes the corresponding MQTT message; publication does not prove that the physical lock has applied the PIN. The source contains the server service, database access and application-to-server flow, but data sharing across multiple devices still requires end-to-end verification on a running deployment. A release build has been checked independently on an Android device in local mode.

\newpage

\begin{center}
    \Large{\textbf{PHÂN CÔNG NHIỆM VỤ}}\\
\end{center}

\vspace{0.8cm}

Bảng dưới đây là phân công tóm tắt. Bảng~\ref{tab:phan-cong-prefinal} ở Chương~\ref{chapter:Assignment} tách từng hạng mục thành tệp, hàm và chỗ được phép sửa khi bị hỏi. Cả hai bảng bám lịch sử Git trên nhánh \texttt{main} sau khi đã gộp tự duyệt phòng thường, lịch lặp, vắng mặt và bộ sơ đồ thiết kế.

\vspace{0.3cm}

\begin{table}[H]
\centering
\small
\renewcommand{\arraystretch}{1.25}

\begin{tabularx}{\textwidth}{
    >{\raggedright\arraybackslash}p{3.6cm}
    >{\raggedright\arraybackslash}X
}
\toprule
\textbf{Thành viên} & \textbf{Phần mã phụ trách trên bản đã merge} \\
\midrule
Đỗ Thùy Linh &
API Java theo nhóm nghiệp vụ, trao đổi ứng dụng--máy chủ, lọc phòng và quy tắc giao khoảng thời gian dùng chung. \\
Hoàng Văn Tùng &
Khung ứng dụng Android, các module màn hình và nghiệp vụ chính, mã PIN theo phòng, OneIoT và mô phỏng SmartLock. \\
Nguyễn Việt Dũng &
Tự duyệt phòng thường, chính sách duyệt khớp giữa ứng dụng và máy chủ, bộ kiểm thử chất lượng. \\
Nguyễn Đình Phúc &
Đặt phòng lặp theo tuần, ghi nhận vắng mặt, thông báo phía quản trị và bộ sơ đồ thiết kế. \\
\bottomrule
\end{tabularx}
\end{table}

\end{document}
